\begin{examplebox}
\textbf{例 1｜基础代入：角度已知，求两段与比例。}
椭圆 $x^2/25+y^2/9=1$，取右焦点 $F=(4,0)$，
$\theta=60^\circ$，$A$ 沿 $60^\circ$，$B$ 沿相反方向。
由 $a=5,b=3,c=4,e=4/5,\cos\theta=1/2$ 得
\[
 AF=\frac9{5+2}=\frac97,\qquad
 BF=\frac9{5-2}=3,\qquad
 \boxed{\frac{AF}{BF}=\frac37},\qquad
 \boxed{AB=\frac{30}{7}}.
\]
\textbf{坐标核验（真正的点，不是随便画的点）：}
\[
 A=\left(\frac{65}{14},\frac{9\sqrt3}{14}\right),\qquad
 B=\left(\frac52,-\frac{3\sqrt3}{2}\right).
\]
两点代回 $x^2/25+y^2/9$ 都恰好等于 $1$；
$A-F$ 与 $(1/2,\sqrt3/2)$ 同向，$B-F$ 与之反向。
\end{examplebox}

\begin{examplebox}
\textbf{例 2｜给斜率，先找余弦的符号。}
仍取 $a=5,b=3,c=4$。过右焦点直线斜率 $k=1$，
故倾斜角 $\theta=45^\circ$，$\cos\theta=\sqrt2/2$。
\[
 \boxed{\frac{AF}{BF}
 =\frac{5-2\sqrt2}{5+2\sqrt2}
 =\frac{33-20\sqrt2}{17}}.
\]
如果 $k=-1$，此时 $\theta=135^\circ$，余弦\textbf{是负数}，
因此比值变为
\[
 \boxed{\frac{AF}{BF}=\frac{5+2\sqrt2}{5-2\sqrt2}>1}.
\]
\textbf{辨识信号：}相同 $|k|$，斜率符号不同，对应正向端点在左还是右。
\end{examplebox}

\begin{examplebox}
\textbf{例 3｜逆向使用：从比值反求直线。}
椭圆 $x^2/4+y^2/3=1$，$F=(1,0)$，
依原端点方向规定，已知 $AF/BF=2$。
有 $a=2,c=1,e=1/2$，故
\[
 \frac{1-\tfrac12\cos\theta}{1+\tfrac12\cos\theta}=2
 \Longrightarrow\ \cos\theta=-\frac23.
\]
$\theta\in[0,\pi)$，所以 $\sin\theta=\sqrt5/3$，
$k=\tan\theta=-\sqrt5/2$：
\[
 \boxed{\ y=-\frac{\sqrt5}{2}(x-1)\ }.
\]
核验：$AF=9/4$，$BF=9/8$；两交点分别为
\[
 A=\left(-\tfrac12,\tfrac{3\sqrt5}{4}\right),\qquad
 B=\left(\tfrac74,-\tfrac{3\sqrt5}{8}\right),
\]
确实在椭圆上。若只说``两段长度的\emph{无序}比为 $2:1$''，
未规定哪段叫 $AF$，则关于 $x$ 轴对称的斜率
$k=\pm\sqrt5/2$ 都需要考虑。
\end{examplebox}

\begin{examplebox}
\textbf{例 4｜转化为弦长最值：斜焦点弦的上、下界。}
由
\[
 AB=\frac{2ab^2}{a^2-c^2\cos^2\theta},\qquad
 0\le \cos^2\theta\le1,
\]
分母在 $[b^2,a^2]$ 之间，因此
\[
 \boxed{\frac{2b^2}{a}\le AB\le2a}.
\]
当 $\theta=\pi/2$（通径，竖直）时取最小值
$2b^2/a$；当 $\theta=0$（长轴，水平）时取最大值 $2a$。
以 $x^2/25+y^2/9=1$ 为例，$\boxed{18/5\le AB\le10}$。
\end{examplebox}

\begin{examplebox}
\textbf{例 5｜知识迁移：换左焦点，哪些符号改变？}
同一椭圆若改取\textbf{左焦点} $F_1=(-c,0)$，
但仍定义 $A=F_1+r(\cos\theta,\sin\theta)$、
$B=F_1-s(\cos\theta,\sin\theta)$（$r,s>0$），则
\[
 \boxed{AF_1=\frac{b^2}{a-c\cos\theta}},\qquad
 \boxed{BF_1=\frac{b^2}{a+c\cos\theta}},
\]
\[
 \boxed{\frac{AF_1}{BF_1}=\frac{1+e\cos\theta}{1-e\cos\theta}}.
\]
记忆不靠新背一套：\textbf{镜像后准线远近对调，所以正负号对调}。
注意 $AF_1$ 表示点 $A$ 到\emph{左焦点}的距离。
\end{examplebox}
